\chapter[1997 --- Prusiner]{1997: Stanley B. Prusiner}
\label{ch:1997_Prusiner}

\section*{Main Contribution}

\textit{Discovered prions---infectious proteins that cause disease without DNA or RNA---challenging the central dogma of molecular biology and explaining devastating neurodegenerative diseases.}

\section{Official Citation}

for his discovery of Prions --- a new biological principle of infection

\section{Background and Historical Context}

In 1982, Stanley B. Prusiner proposed a radical and controversial hypothesis: that a class of infectious diseases---including scrapie in sheep, kuru in humans, and Creutzfeldt--Jakob disease---was caused not by bacteria, viruses, or other conventional pathogens but by a novel type of infectious agent composed entirely of protein. He called these agents ``prions,'' a term derived from ``proteinaceous infectious particle.'' The idea that a protein alone, without DNA or RNA, could be infectious challenged the central dogma of molecular biology and was met with intense skepticism and even hostility from the scientific community. Over the following two decades, however, Prusiner's hypothesis was validated by an overwhelming body of experimental evidence, and his discovery of prions was recognized with the 1997 Nobel Prize in Physiology or Medicine.

Stanley Ben Prusiner was born on May 24, 1942, in Des Moines, Iowa. He studied chemistry at the University of Pennsylvania and then earned his M.D.\ from the University of Pennsylvania School of Medicine in 1968. After completing his residency in neurology at the University of California, San Francisco (UCSF), Prusiner became intrigued by the mysterious infectious agent that caused scrapie and related diseases. In 1972, he began a systematic investigation of this agent, which would consume the next three decades of his career.

The intellectual context of Prusiner's work was shaped by decades of puzzling observations about the ``unconventional'' nature of the scrapie agent. In the 1960s, research by Tikvah Alper and J.S. Griffith had established that the scrapie agent was unusually resistant to procedures that normally inactivate pathogens---it was resistant to ultraviolet radiation, ionizing radiation, heat, and chemicals that destroy nucleic acids. These observations suggested that the agent might not contain DNA or RNA, but no one was willing to propose that a protein alone could be infectious.

Kuru, a fatal neurodegenerative disease that had devastated the Fore people of Papua New Guinea, provided another puzzle. The disease was transmitted through ritualistic cannibalism, but its clinical and pathological features resembled scrapie rather than any known infectious disease. In 1966, Daniel Carleton Gajdusek demonstrated that kuru could be transmitted to chimpanzees, establishing that it was an infectious disease. However, the nature of the infectious agent remained unknown. Gajdusek received the Nobel Prize in 1976 for this work, but the identity of the agent eluded him.

\section{The Discovery/Contribution}

\subsection{Purification of the Prion Protein}

Prusiner's approach was biochemical: he set out to purify the scrapie agent from infected hamster brains and identify its molecular composition. This was an enormous technical challenge, because the agent was present in minute quantities and was significantly difficult to work with. Over the course of several years, Prusiner and his colleagues developed purification protocols that allowed them to enrich the infectious agent by several thousand-fold.

In 1982, Prusiner published a landmark paper in \textit{Science} in which he reported that the purified infectious fraction consisted primarily of a single protein, which he named PrP (prion protein). He proposed that PrP was both the major component of the infectious agent and the agent itself---that the protein was the infectious particle. This proposal was based on the observation that the infectious activity co-purified with PrP and that treatments that destroyed PrP also destroyed infectivity.

\subsection{PrP\texorpdfstring{$^{Sc}$}{Sc} and PrP\texorpdfstring{$^{C}$}{C}: Two Conformations of the Same Protein}

Prusiner's subsequent work revealed a crucial distinction between two forms of PrP. The normal cellular prion protein, PrP$^{C}$, is a glycoprotein anchored to the cell surface of neurons and other cells. It is rich in $\alpha$-helical structure and is soluble in non-denaturing detergents. In contrast, the infectious form, PrP$^{Sc}$ (where ``Sc'' stands for scrapie), is a conformational variant of the same protein. PrP$^{Sc}$ is rich in $\beta$-sheet structure, is insoluble in detergents, and forms aggregates and amyloid fibrils.

The key insight was that PrP$^{Sc}$ could act as a template, inducing the conversion of PrP$^{C}$ into the PrP$^{Sc}$ conformation. This conformational conversion was the fundamental mechanism of prion propagation: when PrP$^{Sc}$ comes into contact with PrP$^{C}$, it catalyzes the refolding of PrP$^{C}$ into the PrP$^{Sc}$ form, creating new infectious particles. This process is self-perpetuating and can be thought of as a ``chain reaction'' of protein misfolding.

\subsection{Inherited Prion Diseases}

Prusiner's work also explained the existence of inherited prion diseases. Certain mutations in the PRNP gene (which encodes PrP) destabilize the PrP$^{C}$ conformation, making the protein more susceptible to spontaneous conversion to the PrP$^{Sc}$ form. These mutations account for the majority of cases of Creutzfeldt--Jakob disease and are also responsible for fatal familial insomnia and Gerstmann--Str{\"a}ussler--Scheinker syndrome. The existence of inherited prion diseases provided powerful support for the protein-only hypothesis, because it demonstrated that a change in protein sequence---not the introduction of an external pathogen---could initiate the disease process.

\subsection{Mad Cow Disease and BSE}

Prusiner's work gained enormous public attention in the 1990s with the emergence of bovine spongiform encephalopathy (BSE), or ``mad cow disease,'' in the United Kingdom. BSE was a prion disease of cattle that was transmitted to humans through the consumption of contaminated meat, causing a variant form of Creutzfeldt--Jakob disease (vCJD). The BSE crisis confirmed the public health significance of prion diseases and underscored the importance of Prusiner's basic research.

\section{Impact on Modern Medicine}

The discovery of prions has had far-reaching implications for medicine, public health, and our understanding of protein folding.

\subsection{Neurodegenerative Disease}

Prion-like mechanisms of protein misfolding and propagation have been implicated in other neurodegenerative diseases, including Alzheimer's disease (involving amyloid-$\beta$ and tau proteins), Parkinson's disease (involving $\alpha$-synuclein), and amyotrophic lateral sclerosis (involving superoxide dismutase 1). The concept that misfolded proteins can spread from cell to cell, seeding the misfolding of normal proteins along the way, is now a central paradigm in neurodegeneration research. This ``prion-like'' spreading mechanism has opened new avenues for therapeutic intervention, including the development of antibodies and small molecules that block the cell-to-cell transmission of misfolded proteins.

\subsection{Diagnostic Advances}

Prusiner's work has enabled the development of diagnostic tests for prion diseases. The detection of PrP$^{Sc}$ in cerebrospinal fluid, blood, and tissue samples provides a basis for early diagnosis. The development of ultrasensitive detection methods, including the real-time quaking-induced conversion (RT-QuIC) assay, has made it possible to detect minute quantities of PrP$^{Sc}$ in clinical specimens, improving the accuracy and speed of diagnosis.

\subsection{Therapeutic Challenges}

Despite the advances in understanding prion biology, effective treatments for prion diseases remain elusive. The misfolded PrP$^{Sc}$ protein is remarkably resistant to degradation and is difficult to target pharmacologically. However, several therapeutic strategies are under investigation, including compounds that stabilize the PrP$^{C}$ conformation, molecules that inhibit the conversion of PrP$^{C}$ to PrP$^{Sc}$, and immunotherapeutic approaches that target PrP$^{Sc}$ for clearance.

\section{Legacy}

Stanley Prusiner became professor of neurology and director of the Institute for Neurodegenerative Diseases at UCSF, where he continues to study prion diseases and their relationship to other neurodegenerative conditions. He received the Albert Lasker Basic Medical Research Award in 1994, three years before the Nobel Prize. Prusiner's willingness to defend his protein-only hypothesis in the face of intense criticism and skepticism---and his ultimate vindication by experimental evidence---stands as a testament to the courage and persistence that are the hallmarks of great science.

His discovery of prions has fundamentally changed our understanding of infectious disease, protein folding, and neurodegeneration. The prion concept has expanded our view of how diseases can spread and how proteins can behave as agents of pathology, with implications that extend far beyond the rare prion diseases to the most common neurodegenerative conditions affecting millions of people worldwide.


\section{Modern Developments}

Prusiner's prion discovery has led to modern neurodegenerative disease research. Understanding of prion-like mechanisms has informed research on Alzheimer's disease (tau and amyloid spreading), Parkinson's disease (alpha-synuclein), and ALS (TDP-43). Anti-prion drugs are in clinical trials for Creutzfeldt-Jakob disease. Understanding of protein misfolding has led to development of chaperone therapies. CRISPR-based diagnostics detect prions with high sensitivity.


\section{Bibliography}

\begin{thebibliography}{9}
\bibitem{nobel-1997}
Nobel Prize Outreach, ``The Nobel Prize in Physiology or Medicine 1997,''\emph{NobelPrize.org}.\\
\url{https://www.nobelprize.org/prizes/medicine/1997/summary/}

\bibitem{lecture-1997}
Stanley B. Prusiner, ``Nobel Lecture,''\emph{NobelPrize.org}. Winner-authored primary source; downloadable from the official Nobel Prize archive.\\
\url{https://www.nobelprize.org/prizes/medicine/1997/prusiner/lecture/}
\end{thebibliography}
